\section{Mathlib Definitions and Theorems}
\label[appendix]{apx:mathlib}
In this appendix, we list verbatim several existing definitions and theorems from the \texttt{mathlib} library that are used in the proofs listed above.

\subsection{Probability Theory}

\subsubsection{\texttt{Mathlib/Probability/Moments/Basic.lean}}
\label{subsection:lean-listing-moments-basic}
\begin{lstlisting}[language=lean]
/-- Moment-generating function of a real random variable `X`: `fun t => μ[exp(t*X)]`. -/
def mgf (X : Ω → ℝ) (μ : Measure Ω) (t : ℝ) : ℝ :=
  μ[fun ω => exp (t * X ω)]

/-- Cumulant-generating function of a real random variable `X`: `fun t => log μ[exp(t*X)]`. -/
def cgf (X : Ω → ℝ) (μ : Measure Ω) (t : ℝ) : ℝ :=
  log (mgf X μ t)

/-- **Chernoff bound** on the upper tail of a real random variable. -/
theorem measure_ge_le_exp_cgf [IsFiniteMeasure μ] (ε : ℝ) (ht : 0 ≤ t)
    (h_int : Integrable (fun ω => exp (t * X ω)) μ) :
    μ.real {ω | ε ≤ X ω} ≤ exp (-t * ε + cgf X μ t)
\end{lstlisting}

\begin{lstlisting}[language=lean]
@[simp]
theorem cgf_zero [IsZeroOrProbabilityMeasure μ] : cgf X μ 0 = 0

theorem mgf_sum_of_identDistrib
    {X : ι → Ω → ℝ}
    {s : Finset ι} {j : ι}
    (h_meas : ∀ i, Measurable (X i))
    (h_indep : iIndepFun X μ)
    (hident : ∀ i ∈ s, ∀ j ∈ s, IdentDistrib (X i) (X j) μ μ)
    (hj : j ∈ s) (t : ℝ) : mgf (∑ i ∈ s, X i) μ t = mgf (X j) μ t ^ #s

lemma mgf_id_map (hX : AEMeasurable X μ) : mgf id (μ.map X) = mgf X μ
\end{lstlisting}

\subsubsection{\texttt{Mathlib/MeasureTheory/Measure/Tilted.lean}}
\begin{lstlisting}[language=lean]
/-- Exponentially tilted measure. When `x ↦ exp (f x)` is integrable, `μ.tilted f` is the
probability measure with density with respect to `μ` proportional to `exp (f x)`. Otherwise it is 0.
-/
noncomputable
def Measure.tilted (μ : Measure α) (f : α → ℝ) : Measure α :=
  μ.withDensity (fun x ↦ ENNReal.ofReal (exp (f x) / ∫ x, exp (f x) ∂μ))
\end{lstlisting}

\begin{lstlisting}[language=lean]
lemma isProbabilityMeasure_tilted [NeZero μ]
    (hf : Integrable (fun x ↦ exp (f x)) μ) :
    IsProbabilityMeasure (μ.tilted f)

@[simp]
lemma tilted_zero' (μ : Measure α) : μ.tilted 0 = (μ Set.univ)⁻¹ • μ

lemma integrable_tilted_iff {E : Type*} [NormedAddCommGroup E] [NormedSpace ℝ E]
    {f : α → ℝ} (hf : Integrable (fun x ↦ exp (f x)) μ) (g : α → E) :
    Integrable g (μ.tilted f) ↔ Integrable (fun x ↦ exp (f x) • g x) μ
\end{lstlisting}

\newpage
\begin{lstlisting}[language=lean]
lemma integral_tilted (f : α → ℝ) (g : α → E) :
    ∫ x, g x ∂(μ.tilted f) = ∫ x, (exp (f x) / ∫ x, exp (f x) ∂μ) • (g x) ∂μ

lemma integral_exp_tilted (f g : α → ℝ) :
    ∫ x, exp (g x) ∂(μ.tilted f) = (∫ x, exp ((f + g) x) ∂μ) / ∫ x, exp (f x) ∂μ

lemma setIntegral_tilted' (f : α → ℝ) (g : α → E) {s : Set α} (hs : MeasurableSet s) :
    ∫ x in s, g x ∂(μ.tilted f) = ∫ x in s, (exp (f x) / ∫ x, exp (f x) ∂μ) • (g x) ∂μ

lemma lintegral_tilted (f : α → ℝ) (g : α → ℝ≥0∞) :
    ∫⁻ x, g x ∂(μ.tilted f)
      = ∫⁻ x, ENNReal.ofReal (exp (f x) / ∫ x, exp (f x) ∂μ) * (g x) ∂μ

lemma tilted_apply' (μ : Measure α) (f : α → ℝ) {s : Set α} (hs : MeasurableSet s) :
    μ.tilted f s = ∫⁻ a in s, ENNReal.ofReal (exp (f a) / ∫ x, exp (f x) ∂μ) ∂μ
\end{lstlisting}

\subsubsection{\texttt{Mathlib/Probability/Moments/Tilted.lean}}
\begin{lstlisting}[language=lean]
/-- The integral of `X` against the tilted measure `μ.tilted (t * X ·)` is the first derivative of
the cumulant-generating function of `X` at `t`. -/
lemma integral_tilted_mul_self (ht : t ∈ interior (integrableExpSet X μ)) :
    (μ.tilted (t * X ·))[X] = deriv (cgf X μ) t

lemma memLp_tilted_mul (ht : t ∈ interior (integrableExpSet X μ)) (p : ℝ≥0) :
    MemLp X p (μ.tilted (t * X ·))

/-- The variance of `X` under the tilted measure `μ.tilted (t * X ·)` is the second derivative of
the cumulant-generating function of `X` at `t`. -/
lemma variance_tilted_mul (ht : t ∈ interior (integrableExpSet X μ)) :
    Var[X; μ.tilted (t * X ·)] = iteratedDeriv 2 (cgf X μ) t
\end{lstlisting}

\subsubsection{\texttt{Mathlib/Probability/StrongLaw.lean}}
\begin{lstlisting}[language=lean]
/-- **Strong law of large numbers**, almost sure version: if `X n` is a sequence of independent
identically distributed integrable real-valued random variables, then `∑ i ∈ range n, X i / n`
converges almost surely to `𝔼[X 0]`. We give here the strong version, due to Etemadi, that only
requires pairwise independence. Superseded by `strong_law_ae`, which works for random variables
taking values in any Banach space. -/
theorem strong_law_ae_real {Ω : Type*} {m : MeasurableSpace Ω} {μ : Measure Ω}
    (X : ℕ → Ω → ℝ) (hint : Integrable (X 0) μ)
    (hindep : Pairwise ((· ⟂ᵢ[μ] ·) on X))
    (hident : ∀ i, IdentDistrib (X i) (X 0) μ μ) :
    ∀ᵐ ω ∂μ, Tendsto (fun n : ℕ => (∑ i ∈ range n, X i ω) / n) atTop (𝓝 μ[X 0])
\end{lstlisting}

\newpage
\subsubsection{\texttt{Mathlib/Probability/Moments/Variance.lean}}
\begin{lstlisting}[language=lean]
/-- The `ℝ`-valued variance of a real-valued random variable defined by applying `ENNReal.toReal`
to `evariance`. -/
def variance : ℝ := (evariance X μ).toReal

/-- **Chebyshev's inequality**: one can control the deviation probability of a real random variable
from its expectation in terms of the variance. -/
theorem meas_ge_le_variance_div_sq [IsFiniteMeasure μ] {X : Ω → ℝ} (hX : MemLp X 2 μ) {c : ℝ}
    (hc : 0 < c) : μ {ω | c ≤ |X ω - μ[X]|} ≤ ENNReal.ofReal (variance X μ / c ^ 2)

theorem variance_smul (c : ℝ) (X : Ω → ℝ) (μ : Measure Ω) :
    variance (c • X) μ = c ^ 2 * variance X μ

lemma variance_eq_integral (hX : AEMeasurable X μ) : Var[X; μ] = ∫ ω, (X ω - μ[X]) ^ 2 ∂μ
\end{lstlisting}

\subsubsection{\texttt{Mathlib/Probability/Moments/IntegrableExpMul.lean}}
\begin{lstlisting}[language=lean]
/-- The interval of reals `t` for which `exp (t * X)` is integrable. -/
def integrableExpSet (X : Ω → ℝ) (μ : Measure Ω) : Set ℝ :=
  {t | Integrable (fun ω ↦ exp (t * X ω)) μ}
\end{lstlisting}

\subsubsection{\texttt{Mathlib/Probability/Independence/Basic.lean}}
\begin{lstlisting}[language=lean]
/-- A family of functions defined on the same space `Ω` and taking values in possibly different
spaces, each with a measurable space structure, is independent if the family of measurable space
structures they generate on `Ω` is independent. For a function `g` with codomain having measurable
space structure `m`, the generated measurable space structure is `MeasurableSpace.comap g m`. -/
def iIndepFun {_mΩ : MeasurableSpace Ω} {β : ι → Type*} [m : ∀ x : ι, MeasurableSpace (β x)]
    (f : ∀ x : ι, Ω → β x) (μ : Measure Ω := by volume_tac) : Prop :=
  Kernel.iIndepFun f (Kernel.const Unit μ) (Measure.dirac () : Measure Unit)

/-- Two functions are independent if the two measurable space structures they generate are
independent. For a function `f` with codomain having measurable space structure `m`, the generated
measurable space structure is `MeasurableSpace.comap f m`.
We use the notation `f ⟂ᵢ[μ] g` for `IndepFun f g μ` (scoped in `ProbabilityTheory`). -/
def IndepFun {β γ} {_mΩ : MeasurableSpace Ω} [MeasurableSpace β] [MeasurableSpace γ]
    (f : Ω → β) (g : Ω → γ) (μ : Measure Ω := by volume_tac) : Prop :=
  Kernel.IndepFun f g (Kernel.const Unit μ) (Measure.dirac () : Measure Unit)
\end{lstlisting}

\begin{lstlisting}[language=lean]
nonrec lemma iIndepFun.comp {β γ : ι → Type*} {mβ : ∀ i, MeasurableSpace (β i)}
    {mγ : ∀ i, MeasurableSpace (γ i)} {f : ∀ i, Ω → β i}
    (h : iIndepFun f μ) (g : ∀ i, β i → γ i) (hg : ∀ i, Measurable (g i)) :
    iIndepFun (fun i ↦ g i ∘ f i) μ

theorem iIndepFun.indepFun {β : ι → Type*}
    {m : ∀ x, MeasurableSpace (β x)} {f : ∀ i, Ω → β i} (hf_Indep : iIndepFun f μ) {i j : ι}
    (hij : i ≠ j) :
    f i ⟂ᵢ[μ] f j

@[to_additive]
lemma iIndepFun.indepFun_finset_prod_of_notMem (hf_Indep : iIndepFun f μ)
    (hf_meas : ∀ i, Measurable (f i)) {s : Finset ι} {i : ι} (hi : i ∉ s) :
    IndepFun (∏ j ∈ s, f j) (f i) μ
\end{lstlisting}
\newpage

\subsubsection{\texttt{Mathlib/Probability/Independence/Integration.lean}}
\begin{lstlisting}[language=lean]
/-- The product of two independent, integrable, real-valued random variables is integrable. -/
theorem IndepFun.integrable_mul {β : Type*} [MeasurableSpace β] {X Y : Ω → β}
    [NormedDivisionRing β] [BorelSpace β] (hXY : X ⟂ᵢ[μ] Y) (hX : Integrable X μ)
    (hY : Integrable Y μ) : Integrable (X * Y) μ

theorem lintegral_prod_eq_prod_lintegral_of_indepFun {ι : Type*}
    (s : Finset ι) (X : ι → Ω → ℝ≥0∞) (hX : iIndepFun X μ)
    (x_mea : ∀ i, Measurable (X i)) :
    ∫⁻ ω, ∏ i ∈ s, (X i ω) ∂μ = ∏ i ∈ s, ∫⁻ ω, X i ω ∂μ
\end{lstlisting}

\subsubsection{\texttt{Mathlib/Probability/IdentDistrib.lean}}
\begin{lstlisting}[language=lean]
/-- Two functions defined on two (possibly different) measure spaces are identically distributed if
their image measures coincide. This only makes sense when the functions are ae measurable
(as otherwise the image measures are not defined), so we require this as well in the definition. -/
structure IdentDistrib (f : α → γ) (g : β → γ)
    (μ : Measure α := by volume_tac)
    (ν : Measure β := by volume_tac) : Prop where
  aemeasurable_fst : AEMeasurable f μ
  aemeasurable_snd : AEMeasurable g ν
  map_eq : Measure.map f μ = Measure.map g ν
\end{lstlisting}

\begin{lstlisting}[language=lean]
protected theorem IdentDistrib.comp {u : γ → δ} (h : IdentDistrib f g μ ν) (hu : Measurable u) :
    IdentDistrib (u ∘ f) (u ∘ g) μ ν

protected theorem IdentDistrib.symm (h : IdentDistrib f g μ ν) : IdentDistrib g f ν μ

protected theorem IdentDistrib.trans {ρ : Measure δ} {h : δ → γ} (h₁ : IdentDistrib f g μ ν)
    (h₂ : IdentDistrib g h ν ρ) : IdentDistrib f h μ ρ

theorem IdentDistrib.integrable_iff [NormedAddCommGroup γ] [BorelSpace γ]
    (h : IdentDistrib f g μ ν) :
    Integrable f μ ↔ Integrable g ν
\end{lstlisting}

\subsubsection{\texttt{Mathlib/Probability/Moments/MGFAnalytic.lean}}
\begin{lstlisting}[language=lean]
/-- The cumulant-generating function is analytic on the interior of the interval
  `integrableExpSet X μ`. -/
lemma analyticOn_cgf :
    AnalyticOn ℝ (cgf X μ) (interior (integrableExpSet X μ))

lemma deriv_cgf (h : v ∈ interior (integrableExpSet X μ)) :
    deriv (cgf X μ) v =
      μ[fun ω ↦ X ω * exp (v * X ω)] / mgf X μ v

lemma deriv_cgf_zero (h : 0 ∈ interior (integrableExpSet X μ)) :
    deriv (cgf X μ) 0 = μ[X] / μ.real Set.univ
\end{lstlisting}
\newpage

\subsubsection{\texttt{Mathlib/MeasureTheory/Measure/Typeclasses/Probability.lean}}
\begin{lstlisting}[language=lean]
lemma prob_le_one {μ : Measure α} [IsZeroOrProbabilityMeasure μ] {s : Set α} : μ s ≤ 1

theorem prob_add_prob_compl [IsProbabilityMeasure μ] (h : MeasurableSet s) : μ s + μ sᶜ = 1

/-- Note that this is not quite as useful as it looks because the measure takes values in `ℝ≥0∞`.
Thus the subtraction appearing is the truncated subtraction of `ℝ≥0∞`, rather than the
better-behaved subtraction of `ℝ`. -/
theorem prob_compl_eq_one_sub (hs : MeasurableSet s) : μ sᶜ = 1 - μ s

@[simp] lemma probReal_univ : μ.real univ = 1
\end{lstlisting}

\subsubsection{\texttt{Mathlib/MeasureTheory/Measure/ProbabilityMeasure.lean}}
\begin{lstlisting}[language=lean]
/-- Probability measures are defined as the subtype of measures that have the property of being
probability measures (i.e., their total mass is one). -/
def ProbabilityMeasure (Ω : Type*) [MeasurableSpace Ω] : Type _ :=
  { μ : Measure Ω // IsProbabilityMeasure μ }

theorem Measure.isProbabilityMeasure_map {f : α → β} (hf : AEMeasurable f μ) :
    IsProbabilityMeasure (map f μ)
\end{lstlisting}

\subsection{Characteristic Functions and Limit Theorems}

\subsubsection{\texttt{Mathlib/MeasureTheory/Measure/CharacteristicFunction/Basic.lean}}
\begin{lstlisting}[language=lean]
/-- The characteristic function of a measure in an inner product space. -/
noncomputable def charFun [Inner ℝ E] (μ : Measure E) (t : E) : ℂ := ∫ x, exp (⟪x, t⟫ * I) ∂μ

lemma charFun_map_mul_comp {X : Type*} {mX : MeasurableSpace X} {μ : Measure X}
    {f : X → ℝ} (hf : AEMeasurable f μ) (r t : ℝ) :
    charFun (μ.map (fun x ↦ r * (f x))) t = charFun (μ.map f) (r * t)
\end{lstlisting}

\subsubsection{\texttt{Mathlib/Probability/Independence/CharacteristicFunction.lean}}
\begin{lstlisting}[language=lean]
lemma iIndepFun.charFun_map_fun_finset_sum_eq_prod [InnerProductSpace ℝ E]
    (mX : ∀ i ∈ s, AEMeasurable (X i) P) (hX : iIndepFun (s.restrict X) P) :
    charFun (P.map (fun ω ↦ ∑ i ∈ s, X i ω)) = ∏ i ∈ s, charFun (P.map (X i))
\end{lstlisting}

\subsubsection{\texttt{Mathlib/Probability/CentralLimitTheorem.lean}}
\begin{lstlisting}[language=lean]
lemma tendsto_charFun_inv_sqrt_mul_pow {X : Ω → ℝ}
    (hX : AEMeasurable X P) (h0 : P[X] = 0) (h1 : P[X ^ 2] = 1) (t : ℝ) :
    Tendsto (fun (n : ℕ) ↦ (charFun (P.map X) ((√n)⁻¹ * t)) ^ n) atTop (𝓝 (exp (- t ^ 2 / 2)))
\end{lstlisting}

\subsubsection{\texttt{Mathlib/MeasureTheory/Measure/LevyConvergence.lean}}
\begin{lstlisting}[language=lean]
/-- The **Lévy convergence theorem**: the weak convergence of probability measures is equivalent
to the pointwise convergence of their characteristic functions. -/
theorem ProbabilityMeasure.tendsto_iff_tendsto_charFun :
    Tendsto μ atTop (𝓝 μ₀) ↔
      ∀ t : E, Tendsto (fun n ↦ charFun (μ n) t) atTop (𝓝 (charFun μ₀ t))
\end{lstlisting}

\newpage
\subsubsection{\texttt{Mathlib/MeasureTheory/Measure/Portmanteau.lean}}
\begin{lstlisting}[language=lean]
theorem ProbabilityMeasure.tendsto_measure_of_null_frontier_of_tendsto' {Ω ι : Type*}
    {L : Filter ι} [MeasurableSpace Ω] [TopologicalSpace Ω] [OpensMeasurableSpace Ω]
    [HasOuterApproxClosed Ω] {μ : ProbabilityMeasure Ω} {μs : ι → ProbabilityMeasure Ω}
    (μs_lim : Tendsto μs L (𝓝 μ)) {E : Set Ω} (E_nullbdry : (μ : Measure Ω) (frontier E) = 0) :
    Tendsto (fun i ↦ (μs i : Measure Ω) E) L (𝓝 ((μ : Measure Ω) E))
\end{lstlisting}

\subsubsection{\texttt{Mathlib/Probability/Distributions/Gaussian/Real.lean}}
\begin{lstlisting}[language=lean]
/-- A Gaussian distribution on `ℝ` with mean `μ` and variance `v`. -/
noncomputable
def gaussianReal (μ : ℝ) (v : ℝ≥0) : Measure ℝ :=
  if v = 0 then Measure.dirac μ else volume.withDensity (gaussianPDF μ v)

/-- The characteristic function of a Gaussian distribution with mean `μ` and variance `v`
is given by `t ↦ exp (t * μ - v * t ^ 2 / 2)`. -/
theorem charFun_gaussianReal (t : ℝ) :
    charFun (gaussianReal μ v) t = cexp (t * μ * I - v * t ^ 2 / 2)

lemma noAtoms_gaussianReal {μ : ℝ} {v : ℝ≥0} (h : v ≠ 0) : NoAtoms (gaussianReal μ v)

lemma gaussianReal_map_neg : (gaussianReal μ v).map (fun x ↦ -x) = gaussianReal (-μ) v
\end{lstlisting}

\subsection{Measure Theory}

\subsubsection{\texttt{Mathlib/MeasureTheory/Measure/Typeclasses/Probability.lean}}
\begin{lstlisting}[language=lean]
/-- A measure `μ` is called a probability measure if `μ univ = 1`. -/
class IsProbabilityMeasure (μ : Measure α) : Prop where
  measure_univ : μ univ = 1
\end{lstlisting}

\subsubsection{\texttt{Mathlib/MeasureTheory/Measure/Typeclasses/Finite.lean}}
\begin{lstlisting}[language=lean]
theorem measure_ne_top (μ : Measure α) [IsFiniteMeasure μ] (s : Set α) : μ s ≠ ∞
\end{lstlisting}

\subsubsection{\texttt{Mathlib/MeasureTheory/Measure/Typeclasses/NoAtoms.lean}}
\begin{lstlisting}[language=lean]
/-- Measure `μ` *has no atoms* if the measure of each singleton is zero.

NB: Wikipedia assumes that for any measurable set `s` with positive `μ`-measure,
there exists a measurable `t ⊆ s` such that `0 < μ t < μ s`. While this implies `μ {x} = 0`,
the converse is not true. -/
class NoAtoms {m0 : MeasurableSpace α} (μ : Measure α) : Prop where
  measure_singleton : ∀ x, μ {x} = 0

theorem Iio_ae_eq_Iic : Iio a =ᵐ[μ] Iic a
theorem Ioi_ae_eq_Ici : Ioi a =ᵐ[μ] Ici a
\end{lstlisting}

\newpage
\subsubsection{\texttt{Mathlib/MeasureTheory/MeasurableSpace/Defs.lean}}
\begin{lstlisting}[language=lean]
@[class] structure MeasurableSpace (α : Type*) where
  /-- Predicate saying that a given set is measurable. Use `MeasurableSet` in the root namespace
  instead. -/
  MeasurableSet' : Set α → Prop
  /-- The empty set is a measurable set. Use `MeasurableSet.empty` instead. -/
  measurableSet_empty : MeasurableSet' ∅
  /-- The complement of a measurable set is a measurable set. Use `MeasurableSet.compl` instead. -/
  measurableSet_compl : ∀ s, MeasurableSet' s → MeasurableSet' sᶜ
  /-- The union of a sequence of measurable sets is a measurable set. Use a more general
  `MeasurableSet.iUnion` instead. -/
  measurableSet_iUnion : ∀ f : ℕ → Set α, (∀ i, MeasurableSet' (f i)) → MeasurableSet' (⋃ i, f i)

/-- `MeasurableSet s` means that `s` is measurable (in the ambient measure space on `α`) -/
def MeasurableSet [MeasurableSpace α] (s : Set α) : Prop :=
  ‹MeasurableSpace α›.MeasurableSet' s

/-- A function `f` between measurable spaces is measurable if the preimage of every
  measurable set is measurable. -/
@[fun_prop]
def Measurable [MeasurableSpace α] [MeasurableSpace β] (f : α → β) : Prop :=
  ∀ ⦃t : Set β⦄, MeasurableSet t → MeasurableSet (f ⁻¹' t)

@[simp, fun_prop]
theorem measurable_const {_ : MeasurableSpace α} {_ : MeasurableSpace β} {a : α} :
    Measurable fun _ : β => a

theorem measurable_id {_ : MeasurableSpace α} : Measurable (@id α)

@[measurability]
protected theorem MeasurableSet.iUnion [Countable ι] ⦃f : ι → Set α⦄
    (h : ∀ b, MeasurableSet (f b)) : MeasurableSet (⋃ b, f b)
@[measurability]
theorem MeasurableSet.iInter [Countable ι] {f : ι → Set α} (h : ∀ b, MeasurableSet (f b)) :
    MeasurableSet (⋂ b, f b)
\end{lstlisting}

\subsubsection{\texttt{Mathlib/MeasureTheory/Group/Arithmetic.lean}}
\begin{lstlisting}[language=lean]
@[to_additive (attr := fun_prop)]
theorem Measurable.mul [MeasurableMul₂ M] (hf : Measurable f) (hg : Measurable g) :
    Measurable fun a => f a * g a
@[to_additive (attr := fun_prop)]
theorem Measurable.const_mul [MeasurableMul M] (hf : Measurable f) (c : M) :
    Measurable fun x => c * f x
@[to_additive (attr := fun_prop)]
theorem Measurable.div_const [MeasurableDiv G] (hf : Measurable f) (c : G) :
    Measurable fun x => f x / c

@[to_additive (attr := measurability, fun_prop)]
theorem Finset.measurable_prod (s : Finset ι) (hf : ∀ i ∈ s, Measurable (f i)) :
    Measurable fun a ↦ ∏ i ∈ s, f i a
\end{lstlisting}

\newpage
\subsubsection{\texttt{Mathlib/MeasureTheory/Function/SpecialFunctions/Basic.lean}}
\begin{lstlisting}[language=lean]
@[fun_prop]
protected theorem Measurable.exp : Measurable fun x => Real.exp (f x)
\end{lstlisting}

\subsubsection{\texttt{Mathlib/MeasureTheory/Constructions/BorelSpace/Order.lean}}
\begin{lstlisting}[language=lean]
@[simp, measurability]
theorem measurableSet_Icc [OrderClosedTopology α] : MeasurableSet (Icc a b)

theorem measurableSet_lt [SecondCountableTopology α] [OrderClosedTopology α]
    {f g : δ → α} (hf : Measurable f) (hg : Measurable g) :
    MeasurableSet { a | f a < g a }

theorem measurableSet_le {f g : δ → α} (hf : Measurable f) (hg : Measurable g) :
    MeasurableSet { a | f a ≤ g a }
\end{lstlisting}

\subsubsection{\texttt{Mathlib/MeasureTheory/Function/L1Space/Integrable.lean}}
\begin{lstlisting}[language=lean]
/-- `Integrable f μ` means that `f` is measurable and that the integral `∫⁻ a, ‖f a‖ ∂μ` is finite.
  `Integrable f` means `Integrable f volume`. -/
@[fun_prop]
def Integrable {α} {_ : MeasurableSpace α} (f : α → ε)
    (μ : Measure α := by volume_tac) : Prop :=
  AEStronglyMeasurable f μ ∧ HasFiniteIntegral f μ

@[fun_prop]
theorem integrable_const [IsFiniteMeasure μ] (c : β) : Integrable (fun _ : α => c) μ

theorem integrable_map_measure {f : α → α'} {g : α' → ε}
    (hg : AEStronglyMeasurable g (Measure.map f μ)) (hf : AEMeasurable f μ) :
    Integrable g (Measure.map f μ) ↔ Integrable (g ∘ f) μ

@[fun_prop]
theorem Integrable.const_mul {f : α → 𝕜} (h : Integrable f μ) (c : 𝕜) :
    Integrable (fun x => c * f x) μ
@[fun_prop]
theorem Integrable.div_const {f : α → 𝕜} (h : Integrable f μ) (c : 𝕜) :
    Integrable (fun x => f x / c) μ
@[fun_prop]
theorem Integrable.sub {f g : α → β} (hf : Integrable f μ) (hg : Integrable g μ) :
    Integrable (f - g) μ

theorem MemLp.integrable {q : ℝ≥0∞} (hq1 : 1 ≤ q) {f : α → ε} [IsFiniteMeasure μ]
    (hfq : MemLp f q μ) : Integrable f μ
\end{lstlisting}

\subsubsection{\texttt{Mathlib/MeasureTheory/Integral/IntegrableOn.lean}}
\begin{lstlisting}[language=lean]
/-- A function is `IntegrableOn` a set `s` if it is almost everywhere strongly measurable on `s`
and if the integral of its pointwise norm over `s` is less than infinity. -/
def IntegrableOn (f : α → ε) (s : Set α) (μ : Measure α := by volume_tac) : Prop :=
  Integrable f (μ.restrict s)
\end{lstlisting}

\newpage
\subsubsection{\texttt{Mathlib/MeasureTheory/Function/LpSeminorm/Defs.lean}}
\begin{lstlisting}[language=lean]
/-- The property that `f : α → E` is a.e. strongly measurable and `(∫ ‖f a‖ ^ p ∂μ) ^ (1/p)`
is finite if `p < ∞`, or `essSup ‖f‖ < ∞` if `p = ∞`. -/
def MemLp [TopologicalSpace ε] (f : α → ε) (p : ℝ≥0∞) (μ : Measure α := by volume_tac) : Prop :=
  AEStronglyMeasurable f μ ∧ eLpNorm f p μ < ∞
\end{lstlisting}

\subsubsection{\texttt{Mathlib/MeasureTheory/Function/LpSeminorm/Basic.lean}}
\begin{lstlisting}[language=lean]
theorem memLp_const (c : E) [IsFiniteMeasure μ] : MemLp (fun _ : α => c) p μ

theorem memLp_map_measure_iff (hg : AEStronglyMeasurable g (Measure.map f μ))
    (hf : AEMeasurable f μ) : MemLp g p (Measure.map f μ) ↔ MemLp (g ∘ f) p μ
\end{lstlisting}

\subsubsection{\texttt{Mathlib/MeasureTheory/Integral/Bochner/Basic.lean}}
\begin{lstlisting}[language=lean]
theorem integral_map {β} [MeasurableSpace β] {φ : α → β} (hφ : AEMeasurable φ μ) {f : β → G}
    (hfm : AEStronglyMeasurable f (Measure.map φ μ)) :
    ∫ y, f y ∂Measure.map φ μ = ∫ x, f (φ x) ∂μ

@[simp]
theorem integral_const (c : E) : ∫ _ : α, c ∂μ = μ.real univ • c

theorem integral_sub {f g : α → G} (hf : Integrable f μ) (hg : Integrable g μ) :
    ∫ a, f a - g a ∂μ = ∫ a, f a ∂μ - ∫ a, g a ∂μ

theorem integral_const_mul {L : Type*} [RCLike L] (r : L) (f : α → L) :
    ∫ a, r * f a ∂μ = r * ∫ a, f a ∂μ

theorem integral_div {L : Type*} [RCLike L] (r : L) (f : α → L) :
    ∫ a, f a / r ∂μ = (∫ a, f a ∂μ) / r

lemma integral_exp_pos {μ : Measure α} {f : α → ℝ} [hμ : NeZero μ]
    (hf : Integrable (fun x ↦ Real.exp (f x)) μ) :
    0 < ∫ x, Real.exp (f x) ∂μ

theorem integral_eq_zero_iff_of_nonneg_ae {f : α → ℝ} (hf : 0 ≤ᵐ[μ] f) (hfi : Integrable f μ) :
    ∫ x, f x ∂μ = 0 ↔ f =ᵐ[μ] 0

theorem ofReal_integral_eq_lintegral_ofReal {f : α → ℝ} (hfi : Integrable f μ) (f_nn : 0 ≤ᵐ[μ] f) :
    ENNReal.ofReal (∫ x, f x ∂μ) = ∫⁻ x, ENNReal.ofReal (f x) ∂μ
\end{lstlisting}

\subsubsection{\texttt{Mathlib/MeasureTheory/Integral/Bochner/Set.lean}}
\begin{lstlisting}[language=lean]
theorem setIntegral_mono_on (hs : MeasurableSet s) (h : ∀ x ∈ s, f x ≤ g x) :
    ∫ x in s, f x ∂μ ≤ ∫ x in s, g x ∂μ

theorem setIntegral_const [CompleteSpace E] (c : E) : ∫ _ in s, c ∂μ = μ.real s • c
\end{lstlisting}

\newpage
\subsubsection{\texttt{Mathlib/MeasureTheory/Integral/Lebesgue/Basic.lean}}
\begin{lstlisting}[language=lean]
theorem Measure.ext_of_lintegral (ν : Measure α)
    (hμν : ∀ f : α → ℝ≥0∞, Measurable f → ∫⁻ a, f a ∂μ = ∫⁻ a, f a ∂ν) : μ = ν
\end{lstlisting}

\subsubsection{\texttt{Mathlib/MeasureTheory/Integral/Lebesgue/Map.lean}}
\begin{lstlisting}[language=lean]
theorem lintegral_map {f : β → ℝ≥0∞} {g : α → β} (hf : Measurable f)
    (hg : Measurable g) : ∫⁻ a, f a ∂map g μ = ∫⁻ a, f (g a) ∂μ
\end{lstlisting}

\subsubsection{\texttt{Mathlib/MeasureTheory/Integral/Lebesgue/Add.lean}}
\begin{lstlisting}[language=lean]
theorem lintegral_mul_const' (r : ℝ≥0∞) (f : α → ℝ≥0∞) (hr : r ≠ ∞) :
    ∫⁻ a, f a * r ∂μ = (∫⁻ a, f a ∂μ) * r
\end{lstlisting}

\subsubsection{\texttt{Mathlib/MeasureTheory/Measure/MeasureSpaceDef.lean}}
\begin{lstlisting}[language=lean]
/-- The real-valued version of a measure. Maps infinite measure sets to zero. Use as `μ.real s`.
The API is developed in `Mathlib/MeasureTheory/Measure/Real.lean`. -/
protected def Measure.real {α : Type*} {m : MeasurableSpace α} (μ : Measure α) (s : Set α) : ℝ :=
  (μ s).toReal
\end{lstlisting}

\subsubsection{\texttt{Mathlib/MeasureTheory/OuterMeasure/Basic.lean}}
\begin{lstlisting}[language=lean]
@[mono, gcongr]
theorem measure_mono (h : s ⊆ t) : μ s ≤ μ t
\end{lstlisting}

\subsubsection{\texttt{Mathlib/MeasureTheory/OuterMeasure/AE.lean}}
\begin{lstlisting}[language=lean]
theorem ae_iff {p : α → Prop} : (∀ᵐ a ∂μ, p a) ↔ μ { a | ¬p a } = 0

theorem ae_of_all {p : α → Prop} (μ : F) : (∀ a, p a) → ∀ᵐ a ∂μ, p a
\end{lstlisting}

\subsubsection{\texttt{Mathlib/MeasureTheory/Measure/MeasureSpace.lean}}
\begin{lstlisting}[language=lean]
/-- Continuity from below: the measure of the union of an increasing sequence of (not necessarily
measurable) sets is the limit of the measures. -/
theorem tendsto_measure_iUnion_atTop [Preorder ι] [IsCountablyGenerated (atTop : Filter ι)]
    {s : ι → Set α} (hm : Monotone s) : Tendsto (μ ∘ s) atTop (𝓝 (μ (⋃ n, s n)))

/-- `μ s + μ sᶜ = μ univ`. -/
theorem measure_add_measure_compl (h : MeasurableSet s) : μ s + μ sᶜ = μ univ

theorem tendsto_measure_Iic_atTop [Preorder α] [(atTop : Filter α).IsCountablyGenerated]
    (μ : Measure α) : Tendsto (fun x => μ (Iic x)) atTop (𝓝 (μ univ))
\end{lstlisting}
\newpage
\subsubsection{\texttt{Mathlib/MeasureTheory/Measure/Map.lean}}
\begin{lstlisting}[language=lean]
@[simp]
theorem map_apply (hf : Measurable f) {s : Set β} (hs : MeasurableSet s) :
    μ.map f s = μ (f ⁻¹' s)

/-- Mapping a measure twice is the same as mapping the measure with the composition. This version is
for measurable functions. See `map_map_of_aemeasurable` when they are just ae measurable. -/
theorem map_map {g : β → γ} {f : α → β} (hg : Measurable g) (hf : Measurable f) :
    (μ.map f).map g = μ.map (g ∘ f)
\end{lstlisting}

\subsection{Analysis}

\subsubsection{\texttt{Mathlib/Analysis/Convex/Integral.lean}}
\begin{lstlisting}[language=lean]
theorem ConvexOn.map_integral_le [IsProbabilityMeasure μ] (hg : ConvexOn ℝ s g)
    (hgc : ContinuousOn g s) (hsc : IsClosed s) (hfs : ∀ᵐ x ∂μ, f x ∈ s) (hfi : Integrable f μ)
    (hgi : Integrable (g ∘ f) μ) : g (∫ x, f x ∂μ) ≤ ∫ x, g (f x) ∂μ

theorem Convex.integral_mem [IsProbabilityMeasure μ] (hs : Convex ℝ s) (hsc : IsClosed s)
    (hf : ∀ᵐ x ∂μ, f x ∈ s) (hfi : Integrable f μ) : (∫ x, f x ∂μ) ∈ s
\end{lstlisting}

\subsubsection{\texttt{Mathlib/Analysis/Convex/Function.lean}}
\begin{lstlisting}[language=lean]
/-- Convexity of functions -/
def ConvexOn : Prop :=
  Convex 𝕜 s ∧ ∀ ⦃x⦄, x ∈ s → ∀ ⦃y⦄, y ∈ s → ∀ ⦃a b : 𝕜⦄, 0 ≤ a → 0 ≤ b → a + b = 1 →
    f (a • x + b • y) ≤ a • f x + b • f y

/-- Strict convexity of functions -/
def StrictConvexOn : Prop :=
  Convex 𝕜 s ∧ ∀ ⦃x⦄, x ∈ s → ∀ ⦃y⦄, y ∈ s → x ≠ y → ∀ ⦃a b : 𝕜⦄, 0 < a → 0 < b → a + b = 1 →
    f (a • x + b • y) < a • f x + b • f y

theorem StrictConvexOn.convexOn {s : Set E} {f : E → β} (hf : StrictConvexOn 𝕜 s f) :
    ConvexOn 𝕜 s f
\end{lstlisting}

\subsubsection{\texttt{Mathlib/Analysis/Convex/SpecificFunctions/Basic.lean}}
\begin{lstlisting}[language=lean]
/-- `Real.exp` is convex on the whole real line. -/
theorem convexOn_exp : ConvexOn ℝ univ exp
\end{lstlisting}

\newpage
\subsubsection{\texttt{Mathlib/Analysis/Convex/Deriv.lean}}
\begin{lstlisting}[language=lean]
/-- Reformulation of `ConvexOn.le_slope_of_hasDerivAt` using `deriv` -/
lemma deriv_le_slope (hfc : ConvexOn ℝ S f) (hx : x ∈ S) (hy : y ∈ S) (hxy : x < y)
    (hfd : DifferentiableAt ℝ f x) :
    deriv f x ≤ slope f x y

/-- Reformulation of `ConvexOn.slope_le_of_hasDerivAt` using `deriv`. -/
lemma slope_le_deriv (hfc : ConvexOn ℝ S f) (hx : x ∈ S) (hy : y ∈ S) (hxy : x < y)
    (hfd : DifferentiableAt ℝ f y) :
    slope f x y ≤ deriv f y

/-- If a function `f` is continuous on a convex set `D ⊆ ℝ` and `f''` is strictly positive on `D`,
then `f` is strictly convex on `D`.
Note that we don't require twice differentiability explicitly as it is already implied by the second
derivative being strictly positive, except at at most one point. -/
theorem strictConvexOn_of_deriv2_pos' {D : Set ℝ} (hD : Convex ℝ D) {f : ℝ → ℝ}
    (hf : ContinuousOn f D) (hf'' : ∀ x ∈ D, 0 < (deriv^[2] f) x) : StrictConvexOn ℝ D f
\end{lstlisting}

\subsubsection{\texttt{Mathlib/Analysis/SpecialFunctions/Exp.lean}}
\begin{lstlisting}[language=lean]
@[continuity]
theorem continuous_exp : Continuous exp
\end{lstlisting}

\subsubsection{\texttt{Mathlib/Analysis/Complex/Exponential.lean}}
\begin{lstlisting}[language=lean]
theorem exp_sum {α : Type*} (s : Finset α) (f : α → ℝ) :
    exp (∑ x ∈ s, f x) = ∏ x ∈ s, exp (f x)

@[bound]
theorem exp_pos (x : ℝ) : 0 < exp x
\end{lstlisting}

\subsubsection{\texttt{Mathlib/Analysis/SpecialFunctions/Log/Basic.lean}}
\begin{lstlisting}[language=lean]
/-- The real logarithm function, equal to the inverse of the exponential for `x > 0`,
to `log |x|` for `x < 0`, and to `0` for `0`. We use this unconventional extension to
`(-∞, 0]` as it gives the formula `log (x * y) = log x + log y` for all nonzero `x` and `y`, and
the derivative of `log` is `1/x` away from `0`. -/
@[pp_nodot]
noncomputable def log (x : ℝ) : ℝ :=
  if hx : x = 0 then 0 else expOrderIso.symm ⟨|x|, abs_pos.2 hx⟩

@[simp, push]
theorem log_exp (x : ℝ) : log (exp x) = x
theorem exp_log (hx : 0 < x) : exp (log x) = x

@[gcongr, bound]
lemma log_le_log (hx : 0 < x) (hxy : x ≤ y) : log x ≤ log y
\end{lstlisting}

\newpage
\subsubsection{\texttt{Mathlib/Data/Real/Sqrt.lean}}
\begin{lstlisting}[language=lean]
/-- The square root of a real number. This returns 0 for negative inputs.

This has notation `√x`. Note that `√x⁻¹` is parsed as `√(x⁻¹)`. -/
@[irreducible] noncomputable def sqrt (x : ℝ) : ℝ :=
  NNReal.sqrt (Real.toNNReal x)

@[simp]
theorem sqrt_pos : 0 < √x ↔ 0 < x
@[simp]
theorem sqrt_mul {x : ℝ} (hx : 0 ≤ x) (y : ℝ) : √(x * y) = √x * √y
@[simp]
theorem sq_sqrt (h : 0 ≤ x) : √x ^ 2 = x
lemma tendsto_sqrt_atTop : Tendsto (√·) atTop atTop
\end{lstlisting}

\subsubsection{\texttt{Mathlib/Analysis/Calculus/Deriv/Basic.lean}}
\begin{lstlisting}[language=lean]
/-- Derivative of `f` at the point `x`, if it exists.  Zero otherwise.

If the derivative exists (i.e., `∃ f', HasDerivAt f f' x`), then
`f x' = f x + (x' - x) • deriv f x + o(x' - x)` where `x'` converges to `x`.
-/
def deriv (f : 𝕜 → F) (x : 𝕜) :=
  fderiv 𝕜 f x 1

/-- `f` has the derivative `f'` at the point `x`.

That is, `f x' = f x + (x' - x) • f' + o(x' - x)` where `x'` converges to `x`.
-/
def HasDerivAt (f : 𝕜 → F) (f' : F) (x : 𝕜) :=
  HasDerivAtFilter f f' (𝓝 x ×ˢ pure x)
\end{lstlisting}

\subsubsection{\texttt{Mathlib/Analysis/Calculus/Deriv/MeanValue.lean}}
\begin{lstlisting}[language=lean]
/-- Let `f : ℝ → ℝ` be a differentiable function. If `f'` is positive, then
`f` is a strictly monotone function.
Note that we don't require differentiability explicitly as it already implied by the derivative
being strictly positive. -/
theorem strictMono_of_deriv_pos {f : ℝ → ℝ} (hf' : ∀ x, 0 < deriv f x) : StrictMono f
\end{lstlisting}

\subsubsection{\texttt{Mathlib/Analysis/Calculus/FDeriv/Defs.lean}}
\begin{lstlisting}[language=lean]
/-- A function `f` is differentiable at a point `x` if it admits a derivative there (possibly
non-unique). -/
@[fun_prop]
def DifferentiableAt (f : E → F) (x : E) :=
  ∃ f' : E →L[𝕜] F, HasFDerivAt f f' x
\end{lstlisting}

\newpage
\subsubsection{\texttt{Mathlib/Analysis/Calculus/ContDiff/Defs.lean}}
\begin{lstlisting}[language=lean]
/-- A function is continuously differentiable up to `n` at a point `x` if, for any integer `k ≤ n`,
there is a neighborhood of `x` where `f` admits derivatives up to order `n`, which are continuous.
The parameter `n` belongs to `WithTop ℕ∞`, i.e., it can be a natural number, `∞`, or `ω`
(when the `ContDiff` scope is open).

For `n = ∞`, we only require that this holds up to any finite order (where the neighborhood may
depend on the finite order we consider).
For `n = ω`, we require the function to be analytic at `x`. The precise
definition we give (all the derivatives should be analytic) is more involved to work around issues
when the space is not complete, but it is equivalent when the space is complete.
-/
@[fun_prop]
def ContDiffAt (n : WithTop ℕ∞) (f : E → F) (x : E) : Prop :=
  ContDiffWithinAt 𝕜 n f univ x
\end{lstlisting}

\subsubsection{\texttt{Mathlib/Analysis/Calculus/IteratedDeriv/Defs.lean}}
\begin{lstlisting}[language=lean]
/-- The `n`-th iterated derivative of a function from `𝕜` to `F`, as a function from `𝕜` to `F`. -/
def iteratedDeriv (n : ℕ) (f : 𝕜 → F) (x : 𝕜) : F :=
  (iteratedFDeriv 𝕜 n f x : (Fin n → 𝕜) → F) fun _ : Fin n => 1

/-- The `n+1`-th iterated derivative can be obtained by differentiating the `n`-th
iterated derivative. -/
theorem iteratedDeriv_succ : iteratedDeriv (n + 1) f = deriv (iteratedDeriv n f)

/-- The `n`-th iterated derivative can be obtained by iterating `n` times the
differentiation operation. -/
theorem iteratedDeriv_eq_iterate : iteratedDeriv n f = deriv^[n] f
\end{lstlisting}

\subsubsection{\texttt{Mathlib/Analysis/Analytic/Basic.lean}}
\begin{lstlisting}[language=lean]
/-- `f` is analytic within `s` if it is analytic within `s` at each point of `s`.  Note that
this is weaker than `AnalyticOnNhd 𝕜 f s`, as `f` is allowed to be arbitrary outside `s`. -/
def AnalyticOn (f : E → F) (s : Set E) : Prop :=
  ∀ x ∈ s, AnalyticWithinAt 𝕜 f s x
\end{lstlisting}

\subsubsection{\texttt{Mathlib/Analysis/Calculus/FDeriv/Analytic.lean}}
\begin{lstlisting}[language=lean]
theorem AnalyticOn.differentiableOn (h : AnalyticOn 𝕜 f s) : DifferentiableOn 𝕜 f s
\end{lstlisting}

\newpage

\subsection{Order and Filters}

\subsubsection{\texttt{Mathlib/Order/Filter/Defs.lean}}
\begin{lstlisting}[language=lean]
/-- `Filter.Tendsto` is the generic "limit of a function" predicate.
  `Tendsto f l₁ l₂` asserts that for every `l₂` neighborhood `a`,
  the `f`-preimage of `a` is an `l₁` neighborhood. -/
def Tendsto (f : α → β) (l₁ : Filter α) (l₂ : Filter β) :=
  l₁.map f ≤ l₂
\end{lstlisting}

\subsubsection{\texttt{Mathlib/Order/LiminfLimsup.lean}}
\begin{lstlisting}[language=lean]
/-- The `limsup` of a function `u` along a filter `f` is the infimum of the `a` such that
the inequality `u x ≤ a` eventually holds for `f`. -/
def limsup (u : β → α) (f : Filter β) : α :=
  limsSup (map u f)

/-- The `liminf` of a function `u` along a filter `f` is the supremum of the `a` such that
the inequality `u x ≥ a` eventually holds for `f`. -/
def liminf (u : β → α) (f : Filter β) : α :=
  limsInf (map u f)

theorem limsup_le_of_le {f : Filter β} {u : β → α} {a}
    (hf : f.IsCoboundedUnder (· ≤ ·) u := by isBoundedDefault)
    (h : ∀ᶠ n in f, u n ≤ a) : limsup u f ≤ a

theorem liminf_le_liminf {α : Type*} [ConditionallyCompleteLattice β] {f : Filter α} {u v : α → β}
    (h : ∀ᶠ a in f, u a ≤ v a)
    (hu : f.IsBoundedUnder (· ≥ ·) u := by isBoundedDefault)
    (hv : f.IsCoboundedUnder (· ≥ ·) v := by isBoundedDefault) :
    liminf u f ≤ liminf v f
\end{lstlisting}

\subsubsection{\texttt{Mathlib/Order/Filter/IsBounded.lean}}
\begin{lstlisting}[language=lean]
lemma isCoboundedUnder_le_of_le [Preorder α] (l : Filter ι) [NeBot l] {f : ι → α} {x : α}
    (hf : ∀ i, x ≤ f i) :
    IsCoboundedUnder (· ≤ ·) l f
\end{lstlisting}

\subsubsection{\texttt{Mathlib/Topology/Order/LiminfLimsup.lean}}
\begin{lstlisting}[language=lean]
/-- If a function has a limit, then its liminf coincides with its limit. -/
theorem Filter.Tendsto.liminf_eq {f : Filter β} {u : β → α} {a : α} [NeBot f]
    (h : Tendsto u f (𝓝 a)) : liminf u f = a
\end{lstlisting}

\newpage
\subsubsection{\texttt{Mathlib/Topology/Order/Basic.lean}}
\begin{lstlisting}[language=lean]
/-- **Squeeze theorem** (also known as **sandwich theorem**). This version assumes that inequalities
hold everywhere. -/
theorem tendsto_of_tendsto_of_tendsto_of_le_of_le [OrderTopology α] {f g h : β → α} {b : Filter β}
    {a : α} (hg : Tendsto g b (𝓝 a)) (hh : Tendsto h b (𝓝 a)) (hgf : g ≤ f) (hfh : f ≤ h) :
    Tendsto f b (𝓝 a)
\end{lstlisting}

\subsubsection{\texttt{Mathlib/Order/Filter/AtTopBot/Defs.lean}}
\begin{lstlisting}[language=lean]
/-- `atTop` is the filter representing the limit `→ ∞` on an ordered set.
It is generated by the collection of up-sets `{b | a ≤ b}`.
(The preorder need not have a top element for this to be well defined,
and indeed is trivial when a top element `x` exists, i.e., it coincides with `pure x`.) -/
@[to_dual
/-- `atBot` is the filter representing the limit `→ -∞` on an ordered set.
It is generated by the collection of down-sets `{b | b ≤ a}`.
(The preorder need not have a bottom element for this to be well defined,
and indeed is trivial when a bottom element `x` exists, i.e., it coincides with `pure x`.) -/]
def atTop [Preorder α] : Filter α :=
  ⨅ a, 𝓟 (Ici a)

@[to_dual eventually_le_atBot]
theorem eventually_ge_atTop [Preorder α] (a : α) : ∀ᶠ x in atTop, a ≤ x
\end{lstlisting}

\subsubsection{\texttt{Mathlib/Order/Filter/AtTopBot/Basic.lean}}
\begin{lstlisting}[language=lean]
@[simp] lemma eventually_atTop : (∀ᶠ x in atTop, p x) ↔ ∃ a, ∀ b ≥ a, p b := mem_atTop_sets
\end{lstlisting}

\subsubsection{\texttt{Mathlib/Order/Filter/AtTopBot/Archimedean.lean}}
\begin{lstlisting}[language=lean]
theorem tendsto_natCast_atTop_atTop [Semiring R] [PartialOrder R] [IsOrderedRing R]
    [Archimedean R] :
    Tendsto ((↑) : ℕ → R) atTop atTop
\end{lstlisting}

\subsubsection{\texttt{Mathlib/Order/Filter/AtTopBot/Field.lean}}
\begin{lstlisting}[language=lean]
/-- If `r` is a positive constant, `fun x ↦ r * f x` tends to infinity along a filter
if and only if `f` tends to infinity along the same filter. -/
theorem tendsto_const_mul_atTop_of_pos (hr : 0 < r) :
    Tendsto (fun x => r * f x) l atTop ↔ Tendsto f l atTop
\end{lstlisting}

\subsubsection{\texttt{Mathlib/Topology/Neighborhoods.lean}}
\begin{lstlisting}[language=lean]
theorem tendsto_const_nhds {f : Filter α} : Tendsto (fun _ : α => x) f (𝓝 x)
\end{lstlisting}

\subsubsection{\texttt{Mathlib/Topology/Order/DenselyOrdered.lean}}
\begin{lstlisting}[language=lean]
@[simp]
theorem frontier_Icc [NoMinOrder α] [NoMaxOrder α] {a b : α} (h : a ≤ b) :
    frontier (Icc a b) = {a, b}
\end{lstlisting}

\subsubsection{\texttt{Mathlib/Order/Bounds/Defs.lean}}
\begin{lstlisting}[language=lean]
/-- A set is bounded above if there exists an upper bound. -/
@[to_dual /-- A set is bounded below if there exists a lower bound. -/]
def BddAbove (s : Set α) :=
  (upperBounds s).Nonempty
\end{lstlisting}

\subsubsection{\texttt{Mathlib/Order/ConditionallyCompleteLattice/Defs.lean}}
\begin{lstlisting}[language=lean]
/-- A conditionally complete lattice is a lattice in which
every nonempty subset which is bounded above has a supremum, and
every nonempty subset which is bounded below has an infimum.
Typical examples are real numbers or natural numbers.

To differentiate the statements from the corresponding statements in (unconditional)
complete lattices, we prefix `sInf` and `sSup` by a `c` everywhere. The same statements should
hold in both worlds, sometimes with additional assumptions of nonemptiness or
boundedness. -/
class ConditionallyCompleteLattice (α : Type*) extends Lattice α, SupSet α, InfSet α where
  /-- Every nonempty subset which is bounded above has a least upper bound. -/
  isLUB_csSup : ∀ s : Set α, s.Nonempty → BddAbove s → IsLUB s (sSup s)
  /-- Every nonempty subset which is bounded below has a greatest lower bound. -/
  isGLB_csInf : ∀ s : Set α, s.Nonempty → BddBelow s → IsGLB s (sInf s)
\end{lstlisting}

\subsubsection{\texttt{Mathlib/Order/ConditionallyCompleteLattice/Indexed.lean}}
\begin{lstlisting}[language=lean]
/-- The indexed supremum of a function is bounded above by a uniform bound -/
theorem ciSup_le [Nonempty ι] {f : ι → α} {c : α} (H : ∀ x, f x ≤ c) : iSup f ≤ c

/-- The indexed supremum of a function is bounded below by the value taken at one point -/
theorem le_ciSup {f : ι → α} (H : BddAbove (range f)) (c : ι) : f c ≤ iSup f
\end{lstlisting}

\subsection{Foundations}

\subsubsection{\texttt{Mathlib/Logic/Pairwise.lean}}
\begin{lstlisting}[language=lean]
/-- A relation `r` holds pairwise if `r i j` for all `i ≠ j`. -/
def Pairwise (r : α → α → Prop) :=
  ∀ ⦃i j⦄, i ≠ j → r i j
\end{lstlisting}

\newpage
\subsection{Extended Reals}

\subsubsection{\texttt{Mathlib/Data/EReal/Basic.lean}}
\begin{lstlisting}[language=lean]
/-- The type of extended real numbers `[-∞, ∞]`, constructed as `WithBot (WithTop ℝ)`. -/
def EReal := WithBot (WithTop ℝ)

@[simp, norm_cast]
theorem coe_mul (x y : ℝ) : (↑(x * y) : EReal) = x * y

@[simp, norm_cast]
protected theorem coe_nonneg {x : ℝ} : (0 : EReal) ≤ x ↔ 0 ≤ x

@[simp]
theorem coe_ne_bot (x : ℝ) : (x : EReal) ≠ ⊥

@[simp]
theorem coe_ne_top (x : ℝ) : (x : EReal) ≠ ⊤
\end{lstlisting}

\subsubsection{\texttt{Mathlib/Data/EReal/Operations.lean}}
\begin{lstlisting}[language=lean]
@[simp, norm_cast] theorem coe_neg (x : ℝ) : (↑(-x) : EReal) = -↑x

lemma mul_nonpos_iff {a b : EReal} : a * b ≤ 0 ↔ 0 ≤ a ∧ b ≤ 0 ∨ a ≤ 0 ∧ 0 ≤ b
\end{lstlisting}

\subsubsection{\texttt{Mathlib/Data/EReal/Inv.lean}}
\begin{lstlisting}[language=lean]
lemma inv_nonneg_of_nonneg {a : EReal} (h : 0 ≤ a) : 0 ≤ a⁻¹
\end{lstlisting}

\subsubsection{\texttt{Mathlib/Topology/Instances/EReal/Lemmas.lean}}
\begin{lstlisting}[language=lean]
protected theorem Tendsto.mul {f : Filter α} {ma : α → EReal} {mb : α → EReal} {a b : EReal}
    (hma : Tendsto ma f (𝓝 a)) (hmb : Tendsto mb f (𝓝 b)) (h₁ : a ≠ 0 ∨ b ≠ ⊥)
    (h₂ : a ≠ 0 ∨ b ≠ ⊤) (h₃ : a ≠ ⊥ ∨ b ≠ 0) (h₄ : a ≠ ⊤ ∨ b ≠ 0) :
    Tendsto (fun x ↦ ma x * mb x) f (𝓝 (a * b))
\end{lstlisting}

\subsubsection{\texttt{Mathlib/Data/ENNReal/Basic.lean}}
\begin{lstlisting}[language=lean]
/-- The extended nonnegative real numbers. This is usually denoted [0, ∞],
  and is relevant as the codomain of a measure. -/
def ENNReal := WithTop ℝ≥0

/-- `toReal x` returns `x` if it is real, `0` otherwise. -/
protected def toReal (a : ℝ≥0∞) : Real := a.toNNReal

/-- `ofReal x` returns `x` if it is nonnegative, `0` otherwise. -/
protected def ofReal (r : Real) : ℝ≥0∞ := r.toNNReal

@[simp]
theorem ofReal_toReal_eq_iff : ENNReal.ofReal a.toReal = a ↔ a ≠ ⊤
\end{lstlisting}

\newpage
\subsubsection{\texttt{Mathlib/Data/ENNReal/Real.lean}}
\begin{lstlisting}[language=lean]
@[gcongr, bound]
theorem ofReal_le_ofReal {p q : ℝ} (h : p ≤ q) : ENNReal.ofReal p ≤ ENNReal.ofReal q

@[gcongr]
theorem toReal_mono (hb : b ≠ ∞) (h : a ≤ b) : a.toReal ≤ b.toReal
\end{lstlisting}

\subsubsection{\texttt{Mathlib/Topology/Instances/ENNReal/Lemmas.lean}}
\begin{lstlisting}[language=lean]
theorem tendsto_toReal {a : ℝ≥0∞} (ha : a ≠ ∞) : Tendsto ENNReal.toReal (𝓝 a) (𝓝 a.toReal)

protected theorem Tendsto.sub {f : Filter α} {ma : α → ℝ≥0∞} {mb : α → ℝ≥0∞} {a b : ℝ≥0∞}
    (hma : Tendsto ma f (𝓝 a)) (hmb : Tendsto mb f (𝓝 b)) (h : a ≠ ∞ ∨ b ≠ ∞) :
    Tendsto (fun a => ma a - mb a) f (𝓝 (a - b))
\end{lstlisting}

\subsubsection{\texttt{Mathlib/Analysis/SpecialFunctions/Log/ENNRealLog.lean}}
\begin{lstlisting}[language=lean]
/-- The logarithm function defined on the extended nonnegative reals `ℝ≥0∞`
to the extended reals `EReal`. Coincides with the usual logarithm function
and with `Real.log` on positive reals, and takes values `log 0 = ⊥` and `log ⊤ = ⊤`.
Conventions about multiplication in `ℝ≥0∞` and addition in `EReal` make the identity
`log (x * y) = log x + log y` unconditional. -/
noncomputable def log (x : ℝ≥0∞) : EReal :=
  if x = 0 then ⊥
    else if x = ⊤ then ⊤
    else Real.log x.toReal

lemma log_ofReal_of_pos {x : ℝ} (hx : 0 < x) : log (ENNReal.ofReal x) = Real.log x

@[gcongr] lemma log_le_log {x y : ℝ≥0∞} (h : x ≤ y) : log x ≤ log y

@[simp] lemma log_le_zero_iff {x : ℝ≥0∞} : log x ≤ 0 ↔ x ≤ 1

@[simp] lemma log_zero : log 0 = ⊥
@[simp] lemma log_one : log 1 = 0
\end{lstlisting}

\subsubsection{\texttt{Mathlib/Analysis/SpecialFunctions/Log/ENNRealLogExp.lean}}
\begin{lstlisting}[language=lean]
@[continuity, fun_prop]
lemma continuous_log : Continuous log
\end{lstlisting}

\subsection{Metric Spaces and Topology}

\subsubsection{\texttt{Mathlib/Topology/MetricSpace/Pseudo/Defs.lean}}
\begin{lstlisting}[language=lean]
theorem tendsto_nhds {f : Filter β} {u : β → α} {a : α} :
    Tendsto u f (𝓝 a) ↔ ∀ ε > 0, ∀ᶠ x in f, dist (u x) a < ε

theorem Real.dist_eq (x y : ℝ) : dist x y = |x - y|
\end{lstlisting}
